\documentclass[10pt,a4paper]{article}
\usepackage[T1]{fontenc}
\usepackage[utf8]{inputenc}
\usepackage[margin=2.2cm]{geometry}
\usepackage{booktabs}
\usepackage{array}
\usepackage{xcolor}
\usepackage{amsmath}
\usepackage[hidelinks]{hyperref}

\definecolor{hdr}{HTML}{1F3864}
\makeatletter
\renewcommand\section{\@startsection{section}{1}{\z@}%
  {-2.4ex \@plus -1ex \@minus -.2ex}{1.3ex \@plus.2ex}%
  {\normalfont\large\bfseries\color{hdr}}}
\renewcommand\subsection{\@startsection{subsection}{2}{\z@}%
  {-2.0ex \@plus -1ex \@minus -.2ex}{1.0ex \@plus.2ex}%
  {\normalfont\normalsize\bfseries\color{hdr}}}
\makeatother
\newenvironment{tightitem}{%
  \begin{itemize}\setlength{\itemsep}{1pt}\setlength{\parskip}{0pt}%
  \setlength{\topsep}{3pt}\setlength{\parsep}{0pt}}{\end{itemize}}
\renewcommand{\arraystretch}{1.15}

\title{\vspace{-1.4cm}\color{hdr}\bfseries Three Ways to Estimate Kinship Under Population Structure\\[2pt]
\large\mdseries\color{black} PC-Relate, PCAngsd and evalAdmix compared,
and how PC-Relate derives $(k_0,k_1,k_2)$}
\author{}
\date{\small Compiled \today}

\begin{document}
\maketitle
\vspace{-1.1cm}

\section{They are the same estimator}

The headline result of this note is that PC-Relate, PCAngsd's kinship estimator and evalAdmix's correlation
of residuals are not three competing ideas. They are one idea, instantiated three times, and they differ only
in bookkeeping. The idea is:

\begin{quote}
Predict each individual's genotype from their \emph{ancestry} alone. Whatever is left over --- the residual
--- is what ancestry cannot explain. Two individuals who share recent ancestors deviate from their
predictions in the \emph{same direction}, so the covariance of their residuals measures relatedness.
\end{quote}

Formally, let $g_{is}\in\{0,1,2\}$ be the allele count for individual $i$ at site $s$, and let
$\pi_{is}$ be the \emph{individual allele frequency} --- the frequency expected at site $s$ given $i$'s
ancestry. Because a genotype is the sum of two alleles,
\begin{equation}
\operatorname{Cov}(g_{is},g_{js}) = 4\phi_{ij}\,\pi_s(1-\pi_s),
\qquad
\operatorname{Var}(g_{is}) = 2\pi_s(1-\pi_s),
\label{eq:cov}
\end{equation}
so the \emph{correlation} of the two residuals is $2\phi_{ij}$ and the covariance normalised by
$4\pi(1-\pi)$ is $\phi_{ij}$ itself. Every method below is a sample version of
Equation~\eqref{eq:cov}. What separates them is exactly three choices:

\begin{tightitem}
  \item[\textbf{(i)}] \textbf{Where $\hat\pi$ comes from} --- a projection onto principal components,
  or a discrete $K$-population admixture model.
  \item[\textbf{(ii)}] \textbf{What goes in the denominator} --- the variance the \emph{model} predicts,
  $\sqrt{\hat\pi(1-\hat\pi)}$, or the variance the residuals \emph{actually have}.
  \item[\textbf{(iii)}] \textbf{Whether the pair's own genotypes were used to build $\hat\pi$} --- and if so,
  whether anything is done about it.
\end{tightitem}

\noindent
Choice (iii) is the one that matters most in practice, and it is where the three methods genuinely diverge.

\section{The three methods}

\subsection{PC-Relate}

$\hat\pi$ comes from a linear regression of genotypes on ancestry-representative principal components: with
$V$ the matrix of PCs (plus an intercept), $2\hat\mu_{is}$ is the fitted value of $g_{is}$.
The residual is $R_{is}=g_{is}-2\hat\mu_{is}$ and
\begin{equation}
\hat\phi^{\text{PC-Relate}}_{ij}
=\frac{\sum_s R_{is}R_{js}}
{4\sum_s \sqrt{\hat\mu_{is}(1-\hat\mu_{is})\,\hat\mu_{js}(1-\hat\mu_{js})}} .
\label{eq:pcrelate}
\end{equation}
The denominator is the \emph{model-expected} residual variance. As instructed, we set aside here the fact
that GENESIS obtains the PCs from PC-AiR on a mutually unrelated training set; that is a property of how the
PCs are built, not of the estimator. PC-Relate is the only one of the three that goes on to produce IBD
sharing probabilities (Section~\ref{sec:ibd}).

\subsection{PCAngsd}

$\hat\pi$ comes from a truncated SVD of the centred genotype \emph{dosages}, iterated: estimate population
frequencies $\hat p_s$, form posterior genotype probabilities and dosages
$E[g_{is}\mid X_{is},\hat\pi_{is}]$ from genotype likelihoods, take the rank-$D$ SVD of the centred dosages,
add $2\hat p_s$ back, and feed $\hat\pi$ in again as the prior until convergence.

The kinship estimator itself is then \emph{identical in form} to
Equation~\eqref{eq:pcrelate}. Reading it off the source of the last version that shipped it
(\texttt{kinship.py}, \texttt{kinship\_cy.pyx}): the numerator matrix is set to $E-2\hat\Pi$, the denominator
matrix to $\sqrt{\hat\Pi(1-\hat\Pi)}$, and the result is their cross-product ratio divided by $4$ ---
Equation~\eqref{eq:pcrelate} with dosages in place of called genotypes. The file is even named for it: its
docstring reads \emph{``Kinship estimator using genotype likelihoods based on PC-Relate''}, and the function
is \texttt{kinshipConomos}. The diagonal is handled separately, from the genotype posteriors, to give
inbreeding coefficients.

So PCAngsd's kinship is PC-Relate, with two substantive changes: $\hat\pi$ is obtained by iterated SVD
rather than one regression, and the input is genotype likelihoods, so nothing is ever hard-called. That
second point is the real contribution --- it is what makes the estimator usable at low depth.

\textbf{Status.} This is no longer available. The \texttt{-kinship} flag, and the companion \texttt{-relate}
option for dropping related individuals from a Beagle file, were removed in the January 2021 Cython rework;
\texttt{kinship.py} and \texttt{kinship\_cy.pyx} were deleted the same day. Current PCAngsd (v1.36) exposes
only \texttt{-{}-inbreed-sites} and \texttt{-{}-inbreed-samples}. PCAone, which re-implements the PCAngsd
algorithm, likewise offers \texttt{-{}-inbreed} but no kinship. For low-depth relatedness today the
supported route is NGSremix, using PCAngsd or NGSadmix as the front-end.

\subsection{evalAdmix}

$\hat\pi$ comes from the admixture model: $\hat\pi_{is}=\sum_k \hat q_{ik}\hat f_{ks}$, with $\hat q$ and
$\hat f$ from ADMIXTURE, STRUCTURE or NGSadmix. The residual is again $r_{is}=g_{is}-2\hat\pi_{is}$, but the
statistic is an ordinary \textbf{Pearson correlation}:
\begin{equation}
c_{ij}=\frac{\sum_s (r_{is}-\bar r_i)(r_{js}-\bar r_j)}
{\sqrt{\sum_s (r_{is}-\bar r_i)^2}\sqrt{\sum_s (r_{js}-\bar r_j)^2}}
\;\longrightarrow\; 2\phi_{ij}.
\label{eq:evaladmix}
\end{equation}
Two differences from Equation~\eqref{eq:pcrelate} follow. First, \textbf{the scale}: a correlation estimates
$2\phi$, not $\phi$, so duplicate pairs return $1$ and full sibs $\tfrac12$. Second, and more interesting,
\textbf{the denominator is empirical} --- the residuals' own observed variance, not
$\hat\pi(1-\hat\pi)$. That makes the statistic self-calibrating: anything that inflates residual variance
uniformly (genotyping error, a slightly wrong $\hat\pi$, unmodelled noise) cancels between numerator and
denominator, where in Equation~\eqref{eq:pcrelate} it would bias $\hat\phi$ downward.

\textbf{The leave-one-out correction.} This is evalAdmix's distinctive feature and the answer to choice
(iii) above. Before correlating individual $i$ with individual $j$, it re-estimates the ancestral frequencies
$F$ with $i$ \emph{excluded} (function \texttt{adaptF}, an EM update run \texttt{-nIts} times, default 5),
recomputes $j$'s residuals under those frequencies, and correlates $i$'s original residual against $j$'s
adapted residual. The point is that if $i$'s own genotypes helped define the frequencies, then $i$'s residual
is shrunk toward zero and the kinship is attenuated. Removing $i$ from the frequency estimation removes that
feedback.

\textbf{The corrected estimator.} Version 1.0 adds \texttt{-method corrected}, the analytic estimator of
van~Waaij et al.\ (2023): instead of leave-one-out refitting, compute the empirical residual correlation
$\hat b$ and subtract the correlation $\hat c$ that the fitted model itself implies, obtained from
$(\mathbf{I}-\mathbf{P})\hat{\mathbf{D}}(\mathbf{I}-\mathbf{P})$ where $\mathbf{P}$ projects onto the column
space of $Q$ (or of the PCs). Single-pass instead of iterative, and it needs only genotypes and $Q$ --- no
\texttt{.P} file. Under a correct model $\hat b-\hat c\to 0$ as sites accumulate.

\textbf{Note on intent.} evalAdmix was written to \emph{test admixture model fit}, not to estimate kinship;
its documentation mentions relatedness only as a caveat (``positive correlation \dots might also be due to
relatedness''). But Equation~\eqref{eq:evaladmix} is a kinship estimator whether or not it was meant as one.

\section{Side by side}

\begin{center}
\small
\begin{tabular}{@{}>{\raggedright\arraybackslash}p{2.6cm}
                  >{\raggedright\arraybackslash}p{3.6cm}
                  >{\raggedright\arraybackslash}p{3.6cm}
                  >{\raggedright\arraybackslash}p{4.4cm}@{}}
\toprule
 & \textbf{PC-Relate} & \textbf{PCAngsd} & \textbf{evalAdmix} \\
\midrule
$\hat\pi$ from & regression on ancestry PCs & iterated truncated SVD on GL dosages & admixture model $\sum_k \hat q_{ik}\hat f_{ks}$ \\
Input & called genotypes & genotype likelihoods & genotypes or GLs (Beagle) \\
Needs $K$? & no & no & yes \\
Statistic & Eq.~\eqref{eq:pcrelate} & Eq.~\eqref{eq:pcrelate} (dosages) & Pearson corr., Eq.~\eqref{eq:evaladmix} \\
Denominator & model, $\sqrt{\hat\pi(1-\hat\pi)}$ & model, $\sqrt{\hat\pi(1-\hat\pi)}$ & empirical residual SD \\
Returns & $\phi$ & $\phi$ & $2\phi$ \\
Self in $\hat\pi$? & yes, uncorrected & yes, uncorrected & \textbf{removed} (leave-one-out or analytic) \\
$(k_0,k_1,k_2)$? & \textbf{yes} & no & no \\
Inbreeding & yes & yes & no \\
Available? & yes (GENESIS) & \textbf{removed in 2021} & yes \\
\bottomrule
\end{tabular}
\end{center}

\section{An empirical comparison}

All three were run on the 126-individual, 104{,}290-SNP example dataset shipped with RelateAdmix ($K=2$),
which is simulated with a known pedigree: ten pairs each of duplicates, parent--offspring, full sibs, half
sibs, first cousins and second cousins, and 7{,}815 unrelated pairs. PC-Relate is the genuine GENESIS
\texttt{pcrelate()}; evalAdmix is the real program, both of its estimators. RelateAdmix is shown as a
reference point --- it is a different kind of method (constrained full ML), not one of the three; see
Section~4.4 before reading its RMSE as a ranking. Exact commands and data
provenance are in Section~7.

\begin{center}
\small
\begin{tabular}{@{}lrrrrrr@{}}
\toprule
\textbf{Relationship} & \textbf{true $\phi$} & \textbf{PC-Relate} & \textbf{PC-Relate} & \textbf{PCAngsd} & \textbf{evalAdmix} & \textbf{(RelateAdmix)} \\
 & & (default) & (corr.\ off) & & & \\
\midrule
unrelated ($n{=}7815$)       & 0      & $+0.00002$ & $-0.0055$ & $-0.0055$ & $-0.0016$ & $+0.0004$ \\
second cousin ($n{=}10$)     & 0.0156 & 0.0200 & 0.0132 & 0.0130 & 0.0177 & 0.0116 \\
first cousin ($n{=}10$)      & 0.0625 & 0.0653 & 0.0564 & 0.0560 & 0.0620 & 0.0526 \\
half sib ($n{=}10$)          & 0.125  & 0.1268 & 0.1199 & 0.1185 & 0.1223 & 0.1108 \\
full sib ($n{=}10$)          & 0.25   & 0.2470 & 0.2410 & 0.2410 & 0.2464 & 0.2439 \\
parent--offspring ($n{=}10$) & 0.25   & 0.2503 & 0.2390 & 0.2374 & 0.2475 & 0.2500 \\
duplicate ($n{=}10$)         & 0.5    & 0.4972 & 0.4860 & 0.4868 & 0.5000 & 0.5000 \\
\midrule
\multicolumn{2}{@{}l}{correlation with truth} & \textbf{0.9891} & 0.9630 & 0.9628 & 0.9883 & 0.9994 \\
\multicolumn{2}{@{}l}{RMSE}                   & \textbf{0.00331} & 0.00824 & 0.00825 & 0.00378 & 0.00089 \\
\multicolumn{2}{@{}l}{SD among unrelated}     & 0.00330 & 0.00616 & 0.00617 & \textbf{0.00345} & 0.00035 \\
\bottomrule
\end{tabular}
\end{center}

\subsection{The whole difference is one post-hoc correction}

The first run of this comparison used a re-implementation of Equation~\eqref{eq:pcrelate} rather than
GENESIS, and produced the ``corr.\ off'' column: attenuation of 3--17\% and unrelated pairs sitting at
$-0.0055$ instead of $0$. The real \texttt{pcrelate()} does not do that. Tracking down why turned out to be
the most informative result here.

It is not the individual allele frequencies. GENESIS computes $\hat\mu = \tfrac12 V(V^{\!\top}V)^{-1}V^{\!\top}G$
(functions \texttt{.calcISAFBeta}, \texttt{.estISAF}) --- an ordinary hat-matrix regression, exactly as
re-implemented. Substituting GENESIS's own PCs changes nothing (unrelated mean $-0.0053$), and its default
\texttt{maf.bound.method = "filter"}, which drops the 20.4\% of sites where some individual's fitted
$\hat\mu$ leaves the MAF range, changes nothing either ($-0.0055$).

The difference is \textbf{\texttt{correctKin}}, enabled by the default \texttt{small.samp.correct = TRUE}.
For each PC $k$, it regresses the kinship estimates on the entries of $V_kV_k^{\!\top}$, fits that line using
\emph{only} pairs with $\hat\phi<2^{-11/2}\approx0.022$ --- i.e.\ the effectively unrelated ones --- and
subtracts it from every pair. Running \texttt{pcrelate()} with \texttt{small.samp.correct = FALSE} reproduces
the re-implementation at a correlation of $\mathbf{1.0000}$, with the unrelated mean identical to five
decimals ($-0.00545$) and the same RMSE ($0.00824$).

So the raw estimator of Equation~\eqref{eq:pcrelate} \emph{is} biased downward, for the reason given under
choice (iii): the pair's own genotypes sit inside $\hat\pi$, the PCs are partly fitted to the relatedness
being measured, and the residuals shrink. PC-Relate does not avoid this --- it measures the bias on the
unrelated pairs and subtracts it afterwards. The correction is worth a factor of $2.5$ in RMSE
($0.00824 \to 0.00331$) and it is what makes PC-Relate competitive.

This also settles the status of PCAngsd's kinship precisely. It implemented
Equation~\eqref{eq:pcrelate} and nothing else --- no \texttt{correctKin}. Its behaviour is therefore the
``corr.\ off'' column, which is why that column and PCAngsd's agree to within $3\times10^{-4}$.
\emph{PCAngsd's kinship was PC-Relate's estimator without PC-Relate's correction.}

\subsection{Four ways to run evalAdmix}

evalAdmix's residual-correlation statistic can be reached by four routes, differing in how the mean structure
is fitted and how the self-contribution is removed. The fourth is the PCAone \texttt{-{}-evaladmix}
implementation of Section~5.1; all four were run on the same data.

\begin{center}
\scriptsize
\begin{tabular}{@{}l cccc@{}}
\toprule
 & \textbf{ADMIXTURE} & \textbf{ADMIXTURE} & \textbf{PCA ($K{-}1$)} & \textbf{PCAone} \\
 & \textbf{+ EM} & \textbf{+ projection} & \textbf{+ projection (R)} & \textbf{\texttt{-{}-evaladmix}} \\
\midrule
Fits mean structure with & $Q,F$ & $Q$ & $K{-}1$ PCs $+$ int. & $K{-}1$ PCs $+$ int. \\
Removes self by & leave-one-out & $\hat b-\hat c$ & $\hat b-\hat c$ & $\hat b-\hat c$ \\
Passes over data & \texttt{-nIts} per ind. & single & single & single \\
Holds genotypes in RAM & yes & yes & yes & \textbf{no} \\
Clips to $[-1,1]$ & --- & no & no & \textbf{yes} \\
\textbf{Needs ADMIXTURE} & \textbf{yes} & \textbf{yes} & \textbf{no} & \textbf{no} \\
\midrule
correlation with truth   & 0.9883 & 0.9883 & 0.9880 & 0.9878 \\
RMSE, related            & 0.00292 & 0.00429 & 0.00443 & \textbf{0.00251} \\
\emph{same, all clipped} & \emph{0.00292} & \emph{0.00255} & \emph{0.00251} & \emph{0.00251} \\
RMSE, related excl.\ dup.& 0.00320 & 0.00279 & \textbf{0.00275} & \textbf{0.00275} \\
SD among unrelated       & \textbf{0.00345} & 0.00349 & 0.00354 & 0.00353 \\
mean, unrelated          & $-0.00156$ & $-0.00154$ & $-0.00155$ & $-0.00155$ \\
mean, duplicates         & \textbf{0.5000} & 0.5079 & 0.5083 & \textbf{0.5000} \\
\midrule
\multicolumn{5}{@{}l}{\emph{attenuation, estimate/truth}} \\
\quad duplicate         & \textbf{1.000} & 1.016 & 1.017 & \textbf{1.000} \\
\quad parent--offspring & 0.990 & 0.998 & \textbf{0.999} & \textbf{0.999} \\
\quad full sib          & 0.985 & 0.990 & \textbf{0.990} & \textbf{0.990} \\
\quad half sib          & 0.979 & 0.981 & \textbf{0.982} & \textbf{0.982} \\
\quad first cousin      & 0.992 & 0.994 & \textbf{0.995} & \textbf{0.995} \\
\quad second cousin     & \textbf{1.134} & 1.135 & 1.136 & 1.136 \\
\midrule
runtime, this dataset & 6\,s wall & 22\,s wall & 2.2\,s & \textbf{0.76\,s wall} \\
                      & 93\,s CPU (20 thr.) & 21\,s CPU (1 thr.) & (R, 1 thr.) & \emph{0.12\,s of which} \\
                      &  &  &  & \emph{is the statistic} \\
\bottomrule
\end{tabular}
\end{center}

\noindent
\textbf{Accuracy is essentially identical across all four} --- correlations with truth differ in the fourth
decimal. The apparent RMSE spread in the first row is almost entirely the $[-1,1]$ clip: apply it uniformly
(second row) and the three projection routes converge on $0.00251$--$0.00255$, with the EM route at
$0.00292$. Excluding duplicates entirely, the projection routes sit at $0.00275$--$0.00279$ against EM's
$0.00320$.

\textbf{Where they genuinely differ is small and systematic.} The EM route is exactly right on duplicates by
construction and is the tightest on unrelated pairs ($\mathrm{SD}=0.00345$); the projection routes are
slightly better on every other class. All four overshoot identically at second cousins ($\approx1.135$),
confirming that is a property of the statistic rather than of any front-end.

\textbf{PCAone \texttt{-{}-evaladmix} combines the two advantages}: projection accuracy on the non-duplicate
classes together with exact duplicates, because it applies the clip that the other projection routes leave
to the user. That gives it the lowest as-shipped RMSE of the four ($0.00251$), though as the clipped row
shows, this is a default rather than a better estimator.

\textbf{The real differences are cost and dependencies.} The EM route refits $F$ with each individual
excluded, \texttt{-nIts} times, so its work grows with sample size on top of the per-pair cost; the
projection routes are single-pass. The two PCA routes need no ADMIXTURE run at all --- invisible here,
because the \texttt{.P}/\texttt{.Q} files ship with the data, but usually the dominant cost of a real
analysis. And only the PCAone implementation avoids holding the genotype matrix in RAM, which is what makes
the approach usable at biobank scale; on this dataset the statistic itself costs $0.12$\,s on top of a PCA
that was going to be run anyway.

\textbf{Recommendation:} \texttt{PCAone -b plink -k <K-1> -{}-evaladmix}. It needs no admixture model, adds
almost nothing to a PCA, scales out-of-core, and matches or beats the alternatives on every relationship
class except the second-cousin overshoot that all four share.

\subsection{Improving the PCA route}

If the PCA $+$ projection route is the one you want, four knobs are worth knowing about. Only two of them
matter, and one matters enormously.

\paragraph{1. Use exactly $K-1$ PCs. This dominates everything else.}
Scanning the number of PCs, with all else fixed:

\begin{center}
\small
\begin{tabular}{@{}lrrrr@{}}
\toprule
PCs used & $k=1$ ($=K{-}1$) & $k=2$ & $k=3$ & $k=4$ \\
\midrule
RMSE (related pairs) & \textbf{0.00443} & 0.03035 & 0.05001 & 0.06917 \\
mean at duplicates ($\phi=0.5$) & \textbf{0.508} & 0.541 & 0.579 & 0.620 \\
correlation with truth & \textbf{0.988} & 0.947 & 0.891 & 0.830 \\
\bottomrule
\end{tabular}
\end{center}

\noindent
One PC too many costs a factor of seven in RMSE. The reason is that with $K=2$ genuinely distinct
populations there is only one real ancestry axis; PC2 onward are fitting \emph{relatedness}, not ancestry, and
projecting them out distorts the residual structure the estimator depends on. What actually has to match is
the \emph{dimension} of the projection: \texttt{method="standard"} appends an intercept row, so $K-1$ PCs
span a $K$-dimensional space, the same as $Q$. Chen \& Storey normalisation (\texttt{method="CS"}) has no
intercept row and therefore needs $k=K$; run at $k=K-1$ it collapses (RMSE $0.0172$, $r=0.67$). Get the
dimension right and the normalisation label barely matters; get it wrong and nothing else can save you.

\paragraph{2. Cap at the parameter bound.} A correlation cannot exceed $1$, so $\phi\le 0.5$; but
$\hat b-\hat c$ is unconstrained and overshoots, returning $0.508$ for duplicates. Truncating at the bound is
as principled as RelateAdmix's simplex constraint and costs nothing:

\begin{center}
\small
\begin{tabular}{@{}lrrr@{}}
\toprule
 & RMSE, related & RMSE, related \emph{excluding} duplicates & duplicates \\
\midrule
PCA $+$ projection            & 0.00443 & 0.00275 & 0.5083 \\
\quad $+$ cap at $\phi=0.5$   & \textbf{0.00251} & 0.00275 & \textbf{0.5000} \\
PCA on unrelated subset       & 0.00415 & 0.00269 & 0.5077 \\
\quad $+$ cap                 & \textbf{0.00246} & \textbf{0.00269} & \textbf{0.5000} \\
\midrule
\emph{evalAdmix, ADMIXTURE $+$ EM} & 0.00292 & 0.00320 & 0.5000 \\
\emph{PC-Relate}                   & 0.00358 & 0.00353 & 0.4972 \\
\bottomrule
\end{tabular}
\end{center}

\noindent
Note the second column carefully. \textbf{Excluding duplicates, the untuned PCA route was already the best
method in this comparison} ($0.00275$, against evalAdmix's $0.00320$ and PC-Relate's $0.00353$). Its entire
apparent deficit in Section~4.2 was ten duplicate pairs breaking a bound. Capping is free and correct, but it
changes nothing unless duplicates are present --- and duplicates are usually removed from real analyses
anyway.

\paragraph{3. Compute the PCs on an unrelated subset and project the rest in.} This is the PC-AiR idea
applied to evalAdmix. The analytic $\hat b-\hat c$ correction handles the \emph{geometry} of the projection,
but not the fact that the PCs themselves were fitted partly to the relatedness being measured. Restricting
the eigendecomposition to a maximal unrelated set (66 of 126 individuals here) and projecting everyone onto
those loadings improves every class slightly. Judged on the 50 related non-duplicate pairs, where the cap
plays no role, RMSE goes from $0.00275$ to $0.00269$ --- \textbf{2.3\%}. That gain is statistically real
(paired $t$ on squared errors $p=0.013$; bootstrap 95\% CI for the RMSE difference
$[2\times10^{-5},\,1.1\times10^{-4}]$, excluding zero; the unrelated-subset version wins in 99.9\% of 10{,}000
bootstrap resamples) but it is practically negligible. And this dataset is an extreme case, with nearly half
the sample sitting in a related pair; expect less in a cohort with few relatives. In practice you would
iterate: run once, identify relatives, recompute the PCs without them. Worth doing if it is cheap; not worth
restructuring a pipeline for.

\paragraph{4. Normalisation and MAF threshold: do not bother.} Centring only versus centring and scaling
versus Chen \& Storey (each at its correct dimension) give RMSE $0.00443$, $0.00441$ and $0.00444$; MAF
thresholds of $0.01$ and $0.05$ are indistinguishable. These are 1\% effects.

\paragraph{What is left --- and the honest summary.} After all of this the PCA route still overshoots at
second cousins by $\approx13\%$, and so do both ADMIXTURE routes by the same amount. That is a property of
the statistic where the signal approaches the noise floor, not of the front-end, and none of these knobs
touches it.

\textbf{So no better approach was found.} Of the four knobs, one (the number of PCs) is a guardrail rather
than an improvement --- it can only make things much worse if set wrong; two (normalisation, MAF) do nothing;
the cap corrects a bound violation that affects duplicates alone; and the unrelated-subset PCA buys 2.3\%.
Run at $K-1$ PCs, \textbf{PCA $+$ projection was already the most accurate method in this entire comparison}
once duplicates are set aside ($0.00275$, against evalAdmix EM's $0.00320$, PC-Relate's $0.00353$ and
RelateAdmix's $0.00812$). The right conclusion is not that it needs improving, but that it was already
winning and the earlier tables obscured that.

\paragraph{Recommended recipe.}
\begin{tightitem}
  \item Choose $K$; use $K-1$ PCs with \texttt{method="standard"} (or $K$ with \texttt{"CS"}).
  \item Run once on all individuals; identify related pairs; recompute the PCs on an unrelated subset and
  project everyone in; re-run.
  \item Truncate the resulting correlations to $[-1,1]$ before halving to the $\phi$ scale.
  \item Do not spend time on normalisation or MAF thresholds.
\end{tightitem}

\subsection{What each method gets right}

Expressed as estimate/truth:

\begin{center}
\small
\begin{tabular}{@{}lrrrrrr@{}}
\toprule
 & duplicate & parent--offs. & full sib & half sib & 1st cousin & 2nd cousin \\
\midrule
PC-Relate (default)   & 0.994 & 1.001 & 0.988 & 1.014 & 1.044 & 1.283 \\
PC-Relate (corr.\ off) / PCAngsd & 0.972 & 0.956 & 0.964 & 0.959 & 0.903 & 0.847 \\
evalAdmix             & 1.000 & 0.990 & 0.985 & 0.979 & 0.992 & 1.134 \\
RelateAdmix           & 1.000 & 1.000 & 0.976 & 0.886 & 0.842 & 0.740 \\
\bottomrule
\end{tabular}
\end{center}

\noindent
PC-Relate and evalAdmix are both close to unbiased down to first cousins and both overshoot at second
cousins, where the signal approaches the noise floor. The uncorrected estimator attenuates monotonically with
distance. RelateAdmix shows the opposite pattern: exact for the closest classes, and attenuating hardest for
the most distant.

\subsection{A warning about overall RMSE --- and about RelateAdmix}

The RMSE column in the main table is misleading, and the way it misleads is worth spelling out.
\textbf{99.24\% of the 7{,}875 pairs are unrelated}, so an RMSE computed over all pairs is very nearly a
measurement of performance on unrelated pairs alone. Splitting it:

\begin{center}
\small
\begin{tabular}{@{}lrrrr@{}}
\toprule
 & \textbf{RMSE, all} & \textbf{RMSE, unrelated} & \textbf{RMSE, related} & \textbf{bias, related} \\
 & (7875 pairs) & (7815) & \textbf{(60)} & \\
\midrule
RelateAdmix              & \textbf{0.00089} & \textbf{0.00054} & 0.00812 & $-0.00571$ \\
PC-Relate                & 0.00331 & 0.00330 & 0.00358 & $+0.00055$ \\
evalAdmix (EM)           & 0.00378 & 0.00378 & 0.00292 & $-0.00120$ \\
PCA $+$ projection       & 0.00387 & 0.00386 & 0.00443 & $+0.00084$ \\
PCAone \texttt{-{}-evaladmix} & 0.00385 & 0.00386 & \textbf{0.00251} & $-0.00055$ \\
PC-Relate (corr.\ off)   & 0.00824 & 0.00823 & 0.00989 & $-0.00842$ \\
\bottomrule
\end{tabular}
\end{center}

\noindent
\textbf{The ranking inverts.} On the 60 pairs that actually have relatedness to estimate, PCAone's
\texttt{-{}-evaladmix} is best ($0.00251$), then evalAdmix's EM route ($0.00292$), then PC-Relate
($0.00358$) --- and \textbf{RelateAdmix is the worst} ($0.00812$), more than three times the best. The two
evalAdmix rows differ only in the $[-1,1]$ clip and the front-end; see Section~4.2. Per class it is the worst method
at half sibs ($0.0152$) and first cousins ($0.0100$), consistent with the attenuation ratios of $0.886$ and
$0.842$ above.

\textbf{Why its unrelated RMSE is so small is a constraint, not precision.} RelateAdmix is a constrained ML
estimator: $(k_0,k_1,k_2)$ must lie on the simplex, so $\phi\ge0$ by construction. Among the 7{,}815
unrelated pairs it never returns a negative value (range $[0,\,0.00384]$) and 93.6\% have
$\hat k_1<10^{-3}$ --- the estimates are pinned at the boundary. PC-Relate, an unconstrained moment-style
estimator, scatters symmetrically about zero (range $[-0.0165,\,+0.0120]$). Pinning at the truth is not the
same as estimating it well; it happens to be right here because the truth for those pairs \emph{is} the
boundary.

The same effect explains RelateAdmix's exactness at the other two boundary classes: duplicates
($k_2=1$) and parent--offspring ($k_1=1$) are simplex vertices, and it returns $0.00000$ RMSE at both.
Its errors are concentrated entirely in the interior --- full sibs, half sibs, cousins --- where no
constraint helps.

\textbf{What this means in practice.} RelateAdmix is excellent at the question ``which pairs are related at
all?'' --- essentially no false positives, and exact on duplicates and parent--offspring. It is the
\emph{least} accurate of these methods at the question ``how related are these two?'' for anything between
full sibs and second cousins. Which one you care about should decide the method, and an overall RMSE will not
tell you.

\subsection{And the IBD probabilities}

PC-Relate's $(k_0,k_1,k_2)$, which neither of the other two produces, recover the pedigree cleanly:

\begin{center}
\small
\begin{tabular}{@{}lrrrrrr@{}}
\toprule
& \multicolumn{3}{c}{\textbf{estimated}} & \multicolumn{3}{c}{\textbf{expected}} \\
\cmidrule(lr){2-4}\cmidrule(lr){5-7}
\textbf{Relationship} & $\hat k_0$ & $\hat k_1$ & $\hat k_2$ & $k_0$ & $k_1$ & $k_2$ \\
\midrule
duplicate         & 0.000 & 0.021 & 0.979 & 0 & 0 & 1 \\
parent--offspring & 0.000 & 1.000 & 0.000 & 0 & 1 & 0 \\
full sib          & 0.255 & 0.502 & 0.243 & 0.25 & 0.50 & 0.25 \\
half sib          & 0.495 & 0.502 & 0.002 & 0.50 & 0.50 & 0 \\
first cousin      & 0.738 & 0.263 & $-0.001$ & 0.75 & 0.25 & 0 \\
\bottomrule
\end{tabular}
\end{center}

\noindent
This is the concrete payoff of Section~6. Parent--offspring and full sibs both have $\phi=1/4$, and indeed
PC-Relate's kinship gives 0.2503 and 0.2470 --- indistinguishable. The IBD probabilities separate them
completely: $k_0=0.000,\,k_1=1.000$ versus $k_0=0.255,\,k_1=0.502,\,k_2=0.243$. No kinship-only method can
do this.

\section{Which one is best?}

\textbf{There is no single ordering, and any claim of one should be treated with suspicion} --- including
the obvious one. PC-Relate has the lower \emph{overall} RMSE ($0.00331$ vs evalAdmix's $0.00378$), but that
is the statistic Section~4.4 warns against: 99.24\% of pairs are unrelated, so it mostly measures
performance on unrelated pairs. Split it and the two swap places depending on the question:

\begin{center}
\small
\begin{tabular}{@{}lrrr@{}}
\toprule
 & \textbf{PC-Relate} & \textbf{evalAdmix (EM)} & \textbf{PCAone \texttt{-{}-evaladmix}} \\
\midrule
RMSE on unrelated pairs (7815) & \textbf{0.00330} & 0.00378 & 0.00386 \\
RMSE on related pairs (60)     & 0.00358 & 0.00292 & \textbf{0.00251} \\
returns $(k_0,k_1,k_2)$        & \textbf{yes} & no & no \\
dimensionality to choose       & \# PCs & $K$ & $K{-}1$ PCs \\
assumes \emph{discrete} ancestry & \textbf{no} & yes & \textbf{no} \\
needs an ADMIXTURE run         & \textbf{no} & yes & \textbf{no} \\
holds genotypes in RAM         & yes & yes & \textbf{no} \\
\bottomrule
\end{tabular}
\end{center}

\noindent
PC-Relate is better at \emph{``are these two unrelated?''}; the two evalAdmix routes are better at
\emph{``how related are they?''}, with the PCAone implementation best of all on related pairs. No margin is
large --- all three correlate with the truth at $0.988$ --- and none is the deciding factor in practice.

What actually separates them is not accuracy at all, but the lower half of that table. PC-Relate is the only
one returning $(k_0,k_1,k_2)$ and inbreeding coefficients. PCAone's \texttt{-{}-evaladmix} is the only one
that needs no admixture run and never holds the genotype matrix in RAM, which is what decides whether the
analysis is possible at all at biobank scale. PCAngsd's kinship no longer exists. These are categorical
differences, and they should drive the choice where a difference in the fourth decimal of an RMSE should
not.

\paragraph{PC-Relate does not escape the dimensionality choice.} It is tempting to say PC-Relate ``needs no
$K$'', and in one sense it does not: PCs do not assume discrete ancestral populations, so continuous
ancestry is handled without pretending otherwise. But the number of PCs entering the regression still has to
be chosen, and that is the same decision wearing different clothes --- with the same sensitivity. Running
\texttt{pcrelate()} on these data with one to four PCs, where $K=2$ makes one correct:

\begin{center}
\small
\begin{tabular}{@{}lrrrr@{}}
\toprule
PCs used & 1 (correct) & 2 & 3 & 4 \\
\midrule
PC-Relate, RMSE on related pairs & \textbf{0.00358} & 0.03474 & 0.06011 & 0.09459 \\
\quad mean at duplicates ($\phi=0.5$) & \textbf{0.4972} & 0.5395 & 0.5880 & 0.6561 \\
\quad correlation with truth & \textbf{0.9891} & 0.9747 & 0.9535 & 0.9319 \\
\midrule
\emph{PCA $+$ projection, for comparison} & \emph{0.00443} & \emph{0.03035} & \emph{0.05001} & \emph{0.06917} \\
\bottomrule
\end{tabular}
\end{center}

\noindent
One PC too many costs a factor of ten, and PC-Relate degrades slightly \emph{faster} than the projection
route it is being compared against. Notably the damage is concentrated in the closest pairs --- duplicates
inflate to $0.66$ while parent--offspring stays at $0.251$ and first cousins at $0.066$ --- so it will not
announce itself in a summary statistic. The real advantage of the PC route is the discreteness assumption it
avoids, not a decision it saves you. A decision rule:

\begin{tightitem}
  \item \textbf{Called genotypes, general use: PC-Relate.} Not because it is the most accurate --- it is
  not, on related pairs --- but because it is complete: inbreeding
  coefficients, and the only IBD probabilities on offer. Keep \texttt{small.samp.correct = TRUE} --- it is
  the default, it is doing most of the work, and it requires all samples in one block.
  \item \textbf{You need to tell parent--offspring from full sibs: PC-Relate, necessarily.} $\phi=1/4$ for
  both; only $k_0$ separates them.
  \item \textbf{evalAdmix} is a close second on $\phi$ (and the best of all on duplicates, exactly $0.5$),
  needs no unrelated set and no leverage correction because leave-one-out removes the problem at source
  rather than after the fact. Use it when you have a defensible $K$, when you want the model-fit diagnostic
  in the same run, or as a cross-check on PC-Relate. Remember it returns $2\phi$.
  \item \textbf{Large samples, or you would rather not run ADMIXTURE: PCAone
  \texttt{-{}-evaladmix}} (Section~5). Same statistic as the evalAdmix projection route, but the mean
  structure comes from $K-1$ PCs, and the single-pass formulation means the genotype matrix never has to fit
  in RAM. It was the most accurate on related pairs here ($0.00251$), it costs $0.12$\,s on top of a PCA you
  were running anyway, and it applies the $\phi\le0.5$ truncation by default. Like the other evalAdmix
  routes it returns $2\phi$ and no IBD probabilities.
  \item \textbf{Ancestry not discrete: PC-Relate or the PCA route}, neither of which assumes discrete
  source populations. This does \emph{not} excuse you from choosing the number of axes --- see above.
  \item \textbf{Low depth: none of the three.} PCAngsd's kinship is gone, and as shown it lacked the
  correction that makes PC-Relate work. Use PCAngsd or NGSadmix for $\hat\pi$ / $\hat q,\hat f$, then
  NGSremix --- ML, genotype likelihoods, and it returns $(k_0,k_1,k_2)$.
  \item \textbf{For screening --- ``is this pair related at all?'' --- RelateAdmix.} Its constrained ML
  never returns a negative $\phi$ and pins unrelated pairs at zero (RMSE $0.00054$ there), so false positives
  are essentially absent, and it is exact on duplicates and parent--offspring. But see Section~4.4: on the
  related pairs themselves it is the \emph{worst} of these methods (RMSE $0.00812$ vs evalAdmix's
  $0.00292$), badly attenuated at half sibs and cousins. Do not read its low overall RMSE as accuracy.
  \item \textbf{Trap:} if you run evalAdmix to choose $K$, close relatives appear as model misfit. Remove
  relatives first, or read the outliers as relatedness.
\end{tightitem}

\noindent
\emph{Scope.} This dataset has discrete, well-separated ancestral populations with $K$ known to be correct,
$n=126$, and no missing data. Small $n$ is exactly the regime \texttt{correctKin} was written for, so its
contribution would be smaller in a biobank; and where ancestry is continuous or $K$ is wrong, $\hat\pi$ from
the admixture model degrades while $\hat\pi$ from PCs does not.

\section{How PC-Relate computes $(k_0,k_1,k_2)$}
\label{sec:ibd}

Since PC-Relate is the only one of the three that does this, it is worth spelling out. What follows is the
GENESIS implementation (\texttt{pcrelate}, \texttt{ibd.probs = TRUE}), read from source.

\paragraph{$k_2$ from a dominance coding.} Define
\begin{equation*}
D_{is}=\begin{cases}\hat\mu_{is} & g_{is}=0\\[1pt] 0 & g_{is}=1\\[1pt] 1-\hat\mu_{is} & g_{is}=2\end{cases}
\qquad
\tilde D_{is}=D_{is}-\hat\mu_{is}(1-\hat\mu_{is}),
\end{equation*}
so that $E[\tilde D_{is}]=0$ under HWE. The dominance covariance between two individuals is proportional to
$k_2$ alone --- it vanishes for IBD-1 sharing --- so
\begin{equation*}
\hat k_2=\frac{\sum_s \tilde D_{is}\tilde D_{js}}
{\sum_s \hat\mu_{is}(1-\hat\mu_{is})\,\hat\mu_{js}(1-\hat\mu_{js})}.
\end{equation*}

\paragraph{$k_0$ from opposing homozygotes.} If a pair shares \emph{any} allele IBD at a site they cannot be
opposing homozygotes. So observed IBS0, over its expectation if the pair shared nothing, estimates $k_0$:
\begin{equation*}
\hat k_0=\frac{\sum_s\left[\mathbf{1}(g_{is}{=}2,g_{js}{=}0)+\mathbf{1}(g_{is}{=}0,g_{js}{=}2)\right]}
{\sum_s\left[\hat\mu_{is}^{2}(1-\hat\mu_{js})^{2}+(1-\hat\mu_{is})^{2}\hat\mu_{js}^{2}\right]}.
\end{equation*}

\paragraph{$k_1$ by subtraction.} $\hat k_1=1-\hat k_0-\hat k_2$. GENESIS stores \texttt{kin}, \texttt{k0}
and \texttt{k2}; $k_1$ is left as the remainder.

\medskip\noindent
\textbf{So it does not go from kinship to IBD probabilities.} $\hat\phi$, $\hat k_0$ and $\hat k_2$ are three
separate statistics computed from the same residuals, and the identity $\phi=k_1/4+k_2/2$ is not used to
derive them. That is deliberate: it is precisely because $\phi$ cannot distinguish parent--offspring from
full sibs that $k_0$ needs its own estimator, and the IBS0 count is what supplies it ($k_0=0$ versus
$k_0=1/4$).

\paragraph{Three corrections, applied only when IBD probabilities are requested.}
\begin{tightitem}
  \item \emph{Inbreeding.} HWE departure inflates the dominance covariance, so the product of the two
  individuals' inbreeding coefficients is subtracted: $\hat k_2\leftarrow\hat k_2-\hat f_i\hat f_j$.
  \item \emph{Residual ancestry and kinship bias in $\hat k_2$.} For each PC $k\ge2$, regress the current
  $\hat k_2$ on the entries of $V_kV_k^{\!\top}$ with linear and quadratic terms, fitting only on pairs with
  $\hat\phi<2^{-11/2}\approx0.022$, and subtract the fit. Then regress on $\hat\phi$ itself, fitting on pairs
  with $\hat k_2<2^{-9/2}\approx0.044$, and subtract that too.
  \item \emph{$k_0$ for distant pairs --- where the kinship identity finally enters.} For pairs below the
  first-degree boundary, $\hat\phi<2^{-5/2}\approx0.177$, the IBS0 estimate is discarded and replaced by
  \begin{equation*}
  \hat k_0=1-4\hat\phi+\hat k_2,
  \end{equation*}
  which is $\phi=k_1/4+k_2/2$ with $k_1=1-k_0-k_2$, rearranged. In distant pairs the IBS0 count is dominated
  by noise and genotyping error while $\hat\phi$ remains well estimated; in first-degree pairs the IBS0
  estimator is kept, for the reason above.
\end{tightitem}

\section{Running this with PCAone}

\subsection{Is there an implementation? No.}

PCAone has no evalAdmix, model-fit or kinship feature --- its only relatedness-adjacent output is
\texttt{-{}-inbreed}. PCAngsd's kinship was removed in 2021. The only existing implementations of the
residual-correlation statistic are the \texttt{evalAdmix} C++ program, which projects onto $Q$ and so needs an
admixture run, and the \texttt{evalPCA()} R code in \texttt{popgenDK/evalPopStructure}, which handles PCA but
holds the whole genotype matrix in memory.

\subsection{But PCAone already computes the expensive part}

PCAone's \texttt{-D/-{}-ld} flag, added for ancestry-adjusted LD, writes a \texttt{.residuals} file that is
precisely $G$ minus its rank-$k$ fit (\texttt{Data.cpp}: \texttt{G -= U * S.asDiagonal() * VT}) --- the
residual matrix the estimator needs. Two things had to be checked, and both hold: the residuals are written
on the \emph{raw} genotype scale, agreeing with residuals computed from unstandardised genotypes to
$1.2\times10^{-7}$ and not with standardised ones (which would differ by $4.5\times10^{-3}$); and the rank-$k$
fit of column-centred data spans the same $K$ dimensions as $k=K-1$ PCs plus an intercept, so
\texttt{-k <K-1>} is the right setting.

\paragraph{Post-hoc recipe.} Run PCAone, then compute $\hat b-\hat c$ from its outputs:
\begin{quote}\ttfamily\small
PCAone -b plink -k <K-1> -d 0 -D -{}-maf 0.05 -o out\\
Rscript evalAdmix\_pcaone.R out plink <K-1>
\end{quote}
Reading \texttt{out.residuals} (\texttt{uint32} $M$, \texttt{uint32} $N$, then $M\times N$ \texttt{float32},
SNP-major) and \texttt{out.eigvecs}:
\begin{equation*}
\hat b = \operatorname{cor}(R),\qquad
\hat c = \operatorname{cov2cor}\!\big((\mathbf{I}-\mathbf{P})\hat{\mathbf{D}}(\mathbf{I}-\mathbf{P})\big),
\qquad \hat\phi = (\hat b-\hat c)/2,
\end{equation*}
with $\mathbf{P}$ the projection onto $[\text{PC}_1\ldots\text{PC}_{K-1},\,\mathbf 1]$ and
$\hat{\mathbf D}$ the diagonal of per-individual mean heterozygosity $\overline{g(2-g)}$. This reproduces
\texttt{evalPCA()} at $r=1.00000$ with a maximum absolute difference of $5.9\times10^{-8}$ --- float
precision. Total cost beyond the PCA: one pass to read the residual file.

\subsection{Implemented: \texttt{-{}-evaladmix} in PCAone}

The algorithm above is implemented in PCAone as \texttt{-{}-evaladmix}
(\texttt{src/EvalAdmix.cpp}, \texttt{src/EvalAdmix.hpp}, wired through \texttt{Cmd} and \texttt{Main}) and is
published at \texttt{github.com/aalbrechtsen/PCAone}, together with the three bug fixes described below:

\begin{quote}\ttfamily\small
PCAone -b plink -k <K-1> -{}-evaladmix -{}-maf 0.05 -o out
\end{quote}

\noindent
It writes \texttt{out.corres} (the $N\times N$ correlation of residuals) and \texttt{out.kinship}
($\text{corres}/2$), both with sample IDs from the \texttt{.fam}. \texttt{-{}-evaladmix-k} selects how many
of the computed PCs enter the projection, so one run can produce many PCs for other purposes while the
statistic uses $K-1$.

\paragraph{Validation.} Against the \texttt{evalPCA()} reference on the RelateAdmix test data:
$r=0.999955$, and identical RMSE on related non-duplicate pairs ($0.00275$ both). The only entries differing
by more than $1.6\times10^{-4}$ are the ten duplicate pairs, where the implementation applies the
$[-1,1]$ clip of Section~4.3 and the R reference does not. In-core and out-of-core runs, and
\texttt{-{}-evaladmix-k 1} against a \texttt{-k 1} run, agree exactly (max difference $0$).
Recovered means: duplicates $0.5000$, parent--offspring $0.2496$, full sibs $0.2475$, half sibs $0.1227$,
first cousins $0.0622$, unrelated $-0.0016$.

\paragraph{Implementation notes.} Three things needed care.
\begin{tightitem}
  \item \emph{Genotype scale.} PCAone codes genotypes as $\{0,\tfrac12,1\}$ (\texttt{BED2GENO}), i.e.\ $g/2$.
  Because $\hat b$ and $\hat c$ are each separately converted to correlations, the constant factor cancels
  and no rescaling is needed.
  \item \emph{Centring.} The out-of-core block reader always returns centred genotypes, ignoring
  \texttt{params.center}. This is harmless for $A$ and $\bar g$ --- per-site centring subtracts $c_s\mathbf 1$
  from column $s$, and every resulting term carries a factor $(\mathbf I-\mathbf P)\mathbf 1=\mathbf 0$
  because the projection contains the intercept --- but it does corrupt the heterozygosity, which is
  therefore recovered as $x = G_{\text{centred}} + \hat f_s$. The in-core path runs with
  \texttt{center = false} and needs no correction.
  \item \emph{The PC scores could not be read correctly.} \texttt{Utils::read\_usv()} was returning a
  transposed matrix for any file with more than one column --- see below. It is now fixed, and this module
  uses it.
\end{tightitem}

\paragraph{Limitations.} PLINK/PGEN input only; \texttt{-{}-maf} with out-of-core is rejected by PCAone
itself, unrelated to this feature; and complete data is assumed, as above. Plain PCA and \texttt{-D} runs are
unaffected.

\subsection{Three bugs found along the way}

Building the feature surfaced three pre-existing defects in PCAone, all now fixed and documented in the
repository's \texttt{docs/} directory. The first is much the most consequential.

\paragraph{1. \texttt{read\_usv()} transposes any multi-column matrix.} The parser fills a flat buffer row by
row, then returned \texttt{Eigen::Map<Mat2D>(V.data(), j, k)} --- and \texttt{Mat2D} is column-major. Mapping
a row-major buffer as column-major scrambles the contents, equivalent to reading the transpose. It is correct
only when there is one column, which is why it went unnoticed.

The victim is not this module but \textbf{the ancestry-adjusted LD path}, \texttt{LD.cpp:493}, which builds
the projector $UU^{\!\top}$ from the result. With $k>1$ the residuals --- and therefore every $R^2$, pruning
and clumping result --- were computed against the wrong subspace. Building both versions and running the
affected path on the test data with $K=2$:

\begin{center}
\small
\begin{tabular}{@{}lr@{}}
\toprule
finite SNP pairs compared & 4{,}611{,}432 \\
$R^2$ values that differ & 4{,}609{,}439 \ (\textbf{99.96\%}) \\
correlation between old and new $R^2$ & \textbf{0.694} \\
mean $R^2$ & $0.01240 \rightarrow 0.01014$ \\
largest single change & $0.2556 \rightarrow 0.0329$ \\
\midrule
pairs with $R^2>0.2$ & 370 (old) $\rightarrow$ 58 (new) \\
flagged by old but not new & 322 \\
\bottomrule
\end{tabular}
\end{center}

\noindent
The old code systematically inflates LD, and at \texttt{-{}-ld-r2 0.2} would prune roughly six times as many
pairs as it should. $R^2$ recomputed independently in R for the six most-affected pairs agrees with the fixed
version exactly in every case, confirming which side is correct.

\paragraph{2. Zero-variance variants give \texttt{NaN} $R^2$.} The LD routines begin from
$1/\texttt{calc\_sds}(G)$. A monomorphic variant, or one lying entirely in the span of the PCs being adjusted
for, has $\mathrm{sd}=0$, so this is $\infty$; the correlation $0\cdot\infty$ is \texttt{NaN}, written to
\texttt{.ld.gz} for every pair involving that variant. \texttt{calc\_sds()} already carried the comment
\emph{``return 1e-9 when sd is 0''} --- the case was anticipated but the guard never written. On the test
data 4 monomorphic variants of 104{,}290 produced 244 \texttt{-nan} rows. Zero-variance columns now yield
$R^2=0$, which is what is true of a variant carrying no information, with a warning naming the count; all
4{,}611{,}432 other values are bit-identical.

\paragraph{3. A missing \texttt{.mbim} segfaults.} \texttt{read\_frq()} opened its file without checking, so
a missing one read as empty, \texttt{F} stayed size zero, and the first \texttt{F(snp\_idx)} indexed out of
bounds --- exit 139, no message. This is easy to reach because \texttt{-P/-{}-USV} points
\texttt{filebim} at \texttt{<prefix>.mbim} and \emph{only} a run with \texttt{-D/-{}-ld} writes one, so a
prefix from a plain PCA run looks valid and crashes. It now names the missing file and exits cleanly.

\subsection{The derivation: why no residual file is needed}

A direct implementation should not write residuals at all. Because $R=G(\mathbf{I}-\mathbf{P})$ and
$\operatorname{cor}$ centres each individual's residuals across sites,
\begin{equation}
\tilde R^{\!\top}\tilde R=(\mathbf{I}-\mathbf{P})\big[G^{\!\top}G-M\,\bar g\,\bar g^{\!\top}\big](\mathbf{I}-\mathbf{P}),
\qquad \hat b=\operatorname{cov2cor}(\tilde R^{\!\top}\tilde R),
\label{eq:gram}
\end{equation}
so the statistic depends on the data only through three quantities:
\begin{tightitem}
  \item $G^{\!\top}G$, the $N\times N$ genotype Gram matrix;
  \item $\bar g$, the per-individual mean genotype;
  \item $\hat{\mathbf D}=\operatorname{diag}\big(\overline{g(2-g)}\big)$, per-individual mean heterozygosity.
\end{tightitem}
All three are simple accumulations over sites, computable in the \emph{same streaming pass} PCAone already
makes, with no second pass, no residual file and no need to hold genotypes. Everything after the pass is
$O(N^3)$ on $N\times N$ matrices. Verified numerically: Equation~\eqref{eq:gram} reproduces the direct
residual computation to $10^{-15}$.

This is a natural fit for PCAone's out-of-core design, and the marginal cost over a PCA it is doing anyway
is one $O(MN^2)$ accumulation --- the same order as forming the covariance matrix it already builds. A
reasonable interface would be a \texttt{-{}-evaladmix} flag emitting the $N\times N$ correlation-of-residuals
matrix, with $k$ defaulting to $K-1$.

\paragraph{Complete data assumed.} Everything in this section assumes no missing genotypes, which is the
case for the data used here. That assumption is what makes the single-pass form exact: with missingness,
evalAdmix skips sites absent in \emph{either} member of a pair, so each pair sees a different set of sites,
$G^{\!\top}G$ ceases to be a sufficient summary, and one would have to accumulate pairwise site counts
alongside it (or mean-impute, at the cost of shrinking the residuals). With complete data no such bookkeeping
is needed and Equation~\eqref{eq:gram} holds exactly.

\paragraph{The algorithm in full.} Given complete genotypes, $K$, and $k=K-1$ PCs:

\begin{quote}\small
\emph{One streaming pass over sites} $s=1\ldots M$, accumulating three objects:\\[2pt]
\hspace*{1em}$A \mathrel{+}= g_s g_s^{\!\top}$ \hfill ($N\times N$; this is $G^{\!\top}G$)\\
\hspace*{1em}$b \mathrel{+}= g_s$ \hfill ($N$; gives $\bar g = b/M$)\\
\hspace*{1em}$d \mathrel{+}= g_s\odot(2-g_s)$ \hfill ($N$; gives $\hat{\mathbf D}=\operatorname{diag}(d/M)$)\\[4pt]
\emph{Then, once:}\\
\hspace*{1em}$\mathbf{P} = V(V^{\!\top}V)^{-1}V^{\!\top}$, \quad $V=[\,\mathrm{PC}_1\ldots\mathrm{PC}_k,\ \mathbf 1\,]$\\
\hspace*{1em}$\hat b = \operatorname{cov2cor}\!\big((\mathbf I-\mathbf P)[A-M\bar g\bar g^{\!\top}](\mathbf I-\mathbf P)\big)$\\
\hspace*{1em}$\hat c = \operatorname{cov2cor}\!\big((\mathbf I-\mathbf P)\hat{\mathbf D}(\mathbf I-\mathbf P)\big)$\\
\hspace*{1em}$\hat\phi = \tfrac12\,\mathrm{clip}(\hat b-\hat c,\,-1,\,1)$
\end{quote}

\noindent
The pass is $O(MN^2)$ --- the same order as the covariance matrix PCAone already forms --- and the tail is
$O(N^3)$. Memory is $O(N^2)$, independent of the number of sites, so nothing about it conflicts with
out-of-core operation. Two points of care remain: $\hat{\mathbf D}$ assumes Hardy--Weinberg within
individuals, so inbred samples need attention; and the $\phi\le0.5$ truncation of Section~4.3 belongs on the
output.

\section{Methods: data and commands}

\subsection{Data}

No data were simulated for this note. The dataset is the example distributed \emph{with} RelateAdmix, in its
\texttt{data/} directory, obtained by cloning the repository:

\begin{quote}\ttfamily\small
git clone https://github.com/aalbrechtsen/relateAdmix.git
\end{quote}

\noindent
It consists of \texttt{smallPlink.\{bed,bim,fam\}} --- 126 individuals, 104{,}290 autosomal SNPs, no missing
genotypes --- together with \texttt{smallPlink.2.P} (ancestral allele frequencies) and
\texttt{smallPlink.2.Q} (admixture proportions), which are ADMIXTURE output for $K=2$ shipped alongside. Because
those two files are provided, ADMIXTURE itself was never run; every method that needs $\hat q,\hat f$ used
the same supplied files, so no method is advantaged by a better admixture fit.

The individuals are simulated from two source populations with a known pedigree. The design is visible in
the data --- individuals are ordered so that relatives are consecutive even--odd pairs --- and was read off
directly: individuals 0--5 unrelated, 6--25 duplicates, 26--45 parent--offspring, 46--65 full sibs, 66--85
half sibs, 86--105 first cousins, 106--125 second cousins, ten pairs per class. ``True $\phi$'' throughout is
the \emph{theoretical} kinship for that relationship ($\tfrac12$, $\tfrac14$, $\tfrac14$, $\tfrac18$,
$\tfrac1{16}$, $\tfrac1{64}$, and $0$ for the remaining 7{,}815 pairs), not any method's output.

\subsection{Software and commands}

\textbf{RelateAdmix} (commit as cloned above) --- C++ command-line build, \texttt{cp CPP\_Makefile Makefile \&\& make}:
\begin{quote}\ttfamily\small
./relateAdmix -plink smallPlink -f smallPlink.2.P -q smallPlink.2.Q -P 20
\end{quote}
5.4\,s wall, 100\,s CPU on 20 threads; output \texttt{output.k} with $(k_0,k_1,k_2)$ per pair.

\medskip\noindent
\textbf{evalAdmix} v1.0, cloned from \texttt{github.com/GenisGE/evalAdmix} and built with \texttt{make}. Both
estimators were run, on the same PLINK files and the same \texttt{.P}/\texttt{.Q}:
\begin{quote}\ttfamily\small
\# default: leave-one-out frequency correction (-nIts 5)\\
./evalAdmix -plink smallPlink -fname smallPlink.2.P -qname smallPlink.2.Q -P 20 -o corres.txt\\[4pt]
\# van Waaij et al.\ analytic projection estimator\\
./evalAdmix -plink smallPlink -fname smallPlink.2.P -qname smallPlink.2.Q -method corrected \textbackslash\\
\hspace*{2em}-P 20 -o corres\_corrected.txt
\end{quote}
Default: 6\,s wall / 93\,s CPU (parallel over 20 threads). Corrected: 22\,s wall / 21.5\,s CPU --- fewer
operations in total, single pass rather than \texttt{-nIts} refits, but effectively single-threaded, so it is
slower in wall time at this sample size and would win as $n$ grows. Both write a $126\times126$ matrix of
correlations of residuals; these were divided by $2$ to put them on the kinship scale. The two agreed closely
(RMSE $0.00378$ vs $0.00382$), so the default estimator is quoted in the main table and the corrected one
reported alongside.

\medskip\noindent
\textbf{evalAdmix, third route (PCA $+$ projection)} --- the R implementation from
\texttt{github.com/popgenDK/evalPopStructure} (\texttt{R/evalPCA.R}), which applies the same $\hat b-\hat c$
estimator to structure from PCA or from the admixture model:
\begin{quote}\ttfamily\small
pca <- makePCA(g, method="standard", center=TRUE, scale=FALSE)\\
ePCA <- evalPCA(pca, k = K-1)\\[4pt]
eADM <- evalPCA(list(method="admixture", g=g, Q=Q), k = K)
\end{quote}
with \texttt{g} the $102{,}185\times126$ SNP-by-individual genotype matrix. The admixture call reproduces the
compiled \texttt{-method corrected} output at $r=1.00000$, confirming the two implementations agree.

\medskip\noindent
\textbf{PC-Relate} --- GENESIS 2.32.0, installed from Bioconductor
(\texttt{BiocManager::install("GENESIS")}, pulling SNPRelate, GWASTools, gdsfmt, SeqArray, SeqVarTools).
The PLINK files were converted with \texttt{SNPRelate::snpgdsBED2GDS()}; ancestry PCs came from
\texttt{snpgdsPCA(maf=0.05, missing.rate=0.05)}; then
\begin{quote}\ttfamily\small
pcrelate(GenotypeBlockIterator(geno, snpBlock=20000),\\
\hspace*{1em}pcs = pcs[,1,drop=FALSE], ibd.probs = TRUE, maf.thresh = 0.05)
\end{quote}
8.5\,s. One PC was used, matching $K=2$. Per the brief, PC-AiR was \emph{not} used --- the PCs come from
plain PCA on all 126 individuals, so no unrelated training set is involved (\texttt{training.set = NULL}).
The run was repeated with \texttt{small.samp.correct = FALSE} to isolate the contribution of
\texttt{correctKin}.

\medskip\noindent
\textbf{PCAngsd} --- \emph{could not be run}, because the kinship estimator no longer exists: the
\texttt{-kinship} flag and \texttt{kinship.py}/\texttt{kinship\_cy.pyx} were removed in January 2021
(commit \texttt{3fda6ed}), and current v1.36 offers only inbreeding. Its formula was therefore read from the
last version that shipped it (\texttt{git show a50c6af\^{}:kinship.py}) and confirmed to be
Equation~\eqref{eq:pcrelate}. The ``PCAngsd'' column is that formula applied to these genotypes with
$\hat\pi$ from a rank-1 SVD of the centred dosages --- which, as Section~4.1 shows, coincides with PC-Relate
with the correction disabled.

\subsection{Code availability}

The \texttt{-{}-evaladmix} implementation and the three fixes above are at
\texttt{github.com/aalbrechtsen/PCAone} (branch \texttt{main}), with
\texttt{docs/evaladmix.md}, \texttt{docs/read-usv-fix.md} and
\texttt{docs/ld-robustness-fixes.md} covering them. A clean clone builds with \texttt{make} and reproduces
every number in Section~4 bit-identically.

\subsection{Reproducibility notes}

Two columns are re-implementations rather than program output, and are labelled as such: ``PCAngsd'' (the
program no longer has the feature) and, in the first version of this comparison, ``PC-Relate''. The latter
was subsequently validated against the genuine \texttt{pcrelate()}: with \texttt{small.samp.correct = FALSE}
the two agree at a correlation of $1.0000$ and identical RMSE, which is what licenses using the
re-implementation to stand in for PCAngsd. All kinship values are on the $\phi$ scale; evalAdmix output was
halved to get there.

\section{References}
\small
\begin{tightitem}
  \item Conomos MP, Miller MB, Thornton TA (2015). Robust inference of population structure for ancestry prediction and correction of stratification in the presence of relatedness. \emph{Genet Epidemiol} 39:276--293. (PC-AiR)
  \item Conomos MP, Reiner AP, Weir BS, Thornton TA (2016). Model-free estimation of recent genetic relatedness. \emph{Am J Hum Genet} 98:127--148. (PC-Relate)
  \item Meisner J, Albrechtsen A (2018). Inferring population structure and admixture proportions in low-depth NGS data. \emph{Genetics} 210:719--731. (PCAngsd)
  \item Garcia-Erill G, Albrechtsen A (2020). Evaluation of model fit of inferred admixture proportions. \emph{Mol Ecol Resour} 20:936--949. (evalAdmix)
  \item van Waaij J, Li S, Garcia-Erill G, Albrechtsen A, Wiuf C (2023). Evaluation of population structure inferred by principal component analysis or the admixture model. \emph{Genetics} 225:iyad157.
  \item Moltke I, Albrechtsen A (2014). RelateAdmix: a software tool for estimating relatedness between admixed individuals. \emph{Bioinformatics} 30:1027--1028.
  \item N\o hr AK, Hangh\o j K, Garcia-Erill G, Li Z, Moltke I, Albrechtsen A (2021). NGSremix: a software tool for estimating pairwise relatedness between admixed individuals from next-generation sequencing data. \emph{G3} 11:jkab174.
  \item Li Z, Meisner J, Albrechtsen A (2023). Fast and accurate out-of-core PCA framework for large scale biobank data. \emph{Genome Research} 33:1599--1608. (PCAone)
\end{tightitem}

\end{document}
